\documentclass[12pt,a4paper]{report}
\usepackage[utf8]{vietnam}
\usepackage[top=2cm, bottom=2cm, left=3cm, right=2cm]{geometry}
\usepackage{graphicx}
\usepackage{float}
\usepackage{tikz}
\usepackage{titlesec}
\usepackage{array}
\usepackage{booktabs}
\usepackage{longtable}
\usepackage{hyperref}
\usepackage[table,xcdraw]{xcolor} % Để tô màu bảng
\usepackage{tabularx}             % Để bảng tự giãn độ rộng
\usepackage{multirow}             % Để gộp dòng
\usepackage{enumitem}             % Để định dạng danh sách trong bảng
\usepackage{longtable} % Thêm gói này vào phần khai báo (Preamble)
\usepackage[table,xcdraw]{xcolor}
\usepackage[table,xcdraw]{xcolor}
\usepackage{xltabular} % Gói quan trọng nhất để ngắt trang và tự căn lề
\usepackage{multirow}
\usetikzlibrary{arrows.meta, positioning, calc}


% Cấu hình tiêu đề chương
\titleformat{\chapter}[display]
  {\normalfont\huge\bfseries}{\chaptername\ \thechapter}{10pt}{\Huge}
\titlespacing*{\chapter}{0pt}{-20pt}{20pt}

\begin{document}

% --- TRANG BÌA ---
\begin{titlepage}
    \begin{center}
        \textbf{ĐẠI HỌC BÁCH KHOA HÀ NỘI}  \\
       \textbf{TRƯỜNG CÔNG NGHỆ THÔNG TIN VÀ TRUYỀN THÔNG}  \\
        \vspace{1cm}
        ****************************************** \\
        \vspace{1.5cm}
        \includegraphics[width=0.4\textwidth]{ảnh latex/hust_logo.png} \\ 
        \vspace{1.5cm}
        {\Huge \textbf{BÁO CÁO GR2}}\\
        \vspace{0.5cm}
        {\Large \textbf{Đề tài: Xây dựng chatbot hỗ trợ học tập}}  \\
        \vfill
        \begin{flushleft}
            \textbf{Giảng viên hướng dẫn:} TS. Trần Đình Khang \\
            \textbf{Sinh viên thực hiện:} Vũ Quang Dũng\\
            \textbf{MSSV:} 20225818 \\
            \textbf{Lớp:} Việt Nhật 01 -K67
        \end{flushleft}
    \end{center}
\end{titlepage}

% --- MỤC LỤC ---
\tableofcontents
\newpage

% --- CHƯƠNG 1 ---
\chapter{Giới thiệu đề tài}
\section{Đặt vấn đề và tổng quan đề tài}
Đối với Đại học Bách khoa Hà Nội, kể từ khi bước sang năm 2 trong chương trình học của mình, các sinh viên bắt đầu tự đăng ký môn học và lớp học cho từng kỳ học và chịu trách nhiệm cho quá trình đăng ký của mình. Trong quá trình đăng ký của sinh viên sẽ gồm hai giai đoạn: đăng ký học phần và đăng ký lớp. 

\textbf{Giai đoạn 1} - đăng ký học phần cho kì tiếp theo sẽ diễn ra từ khá sớm, thường sẽ cách vài tháng trước khi kì học mới diễn ra, ngay sau khi nhà trường có danh sách học phần sẽ mở trong kì. 

\textbf{Giai đoạn 2} - đăng ký lớp sẽ diễn ra trước hi kì học mới bắt đầu khoảng một tháng, sau khi có thời khóa biểu chính thức. Sinh viên sẽ có hai đợt đăng ký trong giai đoạn này. 

\begin{itemize}
    \item \textbf{Đăng ký lớp:} Sinh viên được phép đăng ký các lớp học trong các học phần đã được đăng ký trong giai đoạn 1[cite: 22]. Riêng sinh viên một số chương trình tiên tiến (CTTT) sẽ được đăng ký sớm hơn khoảng 4-7 ngày so với sinh viên các chương trình khác.
    \item \textbf{Đăng ký điều chỉnh:} Sinh viên có thể đăng ký các lớp học trong toàn bộ các học phần mở của kỳ học đó.
\end{itemize}

Sau khi giai đoạn 2 kết thúc, sinh viên hoàn tất quá trình đăng ký của mình. Những sinh viên không đăng ký đủ số tín chỉ tối thiểu của kì học sẽ được các cán bộ phụ trách xếp vào các lớp còn thiếu tương ứng với chương trình đào tạo của sinh viên. 



Hệ thống chatbot hỗ trợ đăng ký học tập sẽ đóng vai trò như một nhà tư vấn, hỗ trợ sinh viên trong cả hai giai đoạn đăng ký. Trước và trong giai đoạn 1, hệ thống đưa ra các phân tích và gợi ý về việc chọn lựa các học phần sao cho phù hợp với chương trình đào tạo, thành tích học tập và nhu cầu cá nhân của sinh viên. Hệ thống cũng có thể đưa ra một số gợi ý về thời khóa biểu các lớp học dựa trên thời khóa biểu tạm thời. Từ khi có thời khóa biểu chính thức và trong giai đoạn 2, hệ thống sẽ giúp sinh viên đưa ra thời khóa biểu phù hợp dựa trên các học phần sinh viên đã đăng ký, dựa trên các sở thích về lịch học, giáo viên hay nhu cầu học cùng bạn bè của sinh viên. 


\section{Mục tiêu và phạm vi đề tài}
Đề tài “Xây dựng hệ thống hỗ trợ đăng ký học tập cho sinh viên” được xây dựng hướng tới mục tiêu hỗ trợ sinh viên đưa ra các phương án đăng ký học phần và đăng ký lớp phù hợp với từng sinh viên, tăng tính cá nhân hóa trong việc hỗ trợ sinh viên đăng ký học tập, nâng cao hiệu quả và năng suất học tập của sinh viên trong suốt kỳ học.
Các mục tiêu cụ thể của đề tài bao gồm: 
\begin{itemize}
    \item Xây dựng hệ thống website dễ sử dụng, giao diện thân thiện, đáp ứng nhu cầu sử dụng nhanh, hiệu quả của sinh viên.
    \item Cá nhân hóa và cụ thể hóa lời khuyên cho sinh viên khi đăng ký học phần và đăng ký lớp.
    \item Hỗ trợ khả năng mở rộng, có thể tích hợp các tính năng khác giúp ích cho hoạt động đăng ký học tập và các hoạt động liên quan đến học tập nói chung cho sinh viên trong tương lai.
\end{itemize}

Phạm vi thực hiện của đề tài tập trung vào phát triển website với các chức năng chính:
\begin{itemize}
    \item Truy vấn, kiểm tra thông tin học tập của từng sinh viên.
    \item Xây dựng thời khóa biểu tạm thời.
    \item Sinh viên trao đổi với chatbot về các vấn đề liên quan đến đăng ký học tập, gợi ý các học phần, lớp học cần đăng ký.
    \item Quản lý tài khoản sinh viên.
    \item Quản lý thông tin học tập của sinh viên.
    \item Quản lý thông tin học phần, lớp mở.
\end{itemize}
Thông qua hệ thống này, đề tài hướng đến khắc phục các hạn chế hiện nay trong việc đăng ký học tập, giúp tăng hiệu quả trong việc học tập của sinh viên.

\section{Định hướng giải pháp}
Để đạt được các mục tiêu trên, đồ án lựa chọn định hướng phát triển hệ thống website hỗ trợ đăng ký học tập dựa trên các công nghệ hiện đại như ReactJS cho giao diện người dùng, FastAPI cho phía backend. Đây là những công nghệ phổ biến, dễ triển khai, đáp ứng tốt yêu cầu về quản lý thông tin tập trung và truy cập đa người dùng.


Giải pháp của đồ án là xây dựng website có đầy đủ chức năng như quản lý học phần, lớp học, thông tin học tập của sinh viên, chatbot trao đổi với sinh viên và xây dựng thời khóa biểu theo mong muốn.


Đóng góp chính của đồ án là phát triển nền tảng hỗ trợ đăng ký học tập hiệu quả, kết quả đạt được là hệ thống có khả năng sử dụng thực tế, dễ dàng mở rộng và tích hợp thêm các chức năng mới nhằm đáp ứng nhu cầu ngày càng tăng cao của sinh viên.


\section{Bố cục đồ án}
Phần còn lại của báo cáo đồ án tốt nghiệp này được tổ chức như sau: 

Chương 2: Khảo sát và phân tích yêu cầu: Trình bày tổng quan các vấn đề liên quan đến việc đăng ký học tập của sinh viên và các chức năng, yêu cầu chức năng của hệ thống.

 Chương 3: Công nghệ sử dụng: Giới thiệu và phân tích các công nghệ, công cụ hiện đại phục vụ phát triển hệ thống, các tiêu chí lựa chọn công nghệ này và phương pháp tiếp cận kiến trúc hệ thống, lý do lựa chọn mô hình phát triển.

 Chương 4: Thiết kế, triển khai và đánh giá: Phân tích và thiết kế hệ thống hỗ trợ đăng ký học tập, đồng thời trình bày các bước triển khai giao diện người dùng, thiết kế cơ sở dữ liệu và tích hợp các chức năng chính.

 Chương 5: Kết luận và hướng phát triển: Trình bày kết quả thực nghiệm, đưa ra kết luận và hướng phát triển, mở rộng trong tương lai. 


% --- CHƯƠNG 2 ---
\chapter{Phân tích đề tài}
\section{Khảo sát hiện trạng}
Trong nhiều năm qua, sinh viên gặp khó khăn trong cả 2 giai đoạn đăng ký học phần và đăng ký lớp. Việc đăng ký học tập của sinh viên đang gặp nhiều khó khăn do nhu cầu cá nhân hóa trong việc lựa chọn môn học và lớp học. Sinh viên phải đối mặt với nhiều vấn đề như:
\begin{itemize}
    \item Trượt môn: Nhiều sinh viên do lơ là, chểnh mảng trong học tập, hoặc vì một số lý do khách quan như sức khỏe, nên bị trượt môn và phải học lại. Những sinh viên này cần học lại môn sớm nhất có thể để tránh bị nâng mức cảnh báo.
    \item Mong muốn học cải thiện: Một số sinh viên tuy không trượt môn nhưng điểm số của một số môn thấp hơn so với mong muốn hay điểm tích lũy CPA còn thiếu một chút để có thể nâng mức xếp loại học lực. Những sinh viên này mong muốn học cải thiện một số môn học nhất định hoặc những môn điểm kém, nhưng không biết có nên học cải thiện hay không và nên học môn nào.
    \item Nhu cầu học vượt: Một số sinh viên muốn tốt nghiệp sớm hoặc học các môn có kiến thức phù hợp với mục tiêu đi làm hoặc thực tập, vì vậy các bạn muốn đăng ký học trước một số môn của kì sau.
    \item Vướng các hoạt động khác: Đời sống của các bạn sinh viên không chỉ có học trên ghế nhà trường mà các bạn còn muốn sắp xếp thời gian học nhiều kỹ năng bên ngoài, tham gia các hoạt động tình nguyện, các bộ môn năng khiếu, thể thao hay tham gia hoạt động của đoàn, trường, câu lạc bộ, kiếm điểm rèn luyện và thực tập, đi làm thêm. Vì vậy, lịch học cần phân bổ hợp lý theo từng nhu cầu cá nhân.
    \item Thời khóa biểu không tối ưu: Nhiều sinh viên khi đăng ký được lớp rồi lại thấy không hài lòng với thời khóa biểu. Một số lý do kể đến như là một ngày học quá nhiều môn, gây quá tải trong ngày học; một ngày có quá ít môn, dẫn đến phải di chuyển nhiều lần từ nhà đến trường; không thích các thầy cô dạy môn đó; không được học cùng bạn, khoảng cách di chuyển giữa các phòng học lớn trong khi thời gian nghỉ ít; thời gian nghỉ giữa các tiết quá ngắn hoặc quá dài; lượng kiến thức trong một kỳ quá nặng hoặc quá nhẹ, số tin chỉ quá ít hoặc quá nhiều, quá ít hoặc quá nhiều môn chuyên ngành, các môn học theo sát chương trình chuẩn, còn nhiều môn nợ chưa được đăng kí, không thích môn thể chất hoặc bổ trợ đã chọn.
    \item Chưa có chức năng xếp thời khóa biểu: Sinh viên vẫn phải truy cập trang web ngoài nhà trường để tự xếp thời khóa biểu thủ công, tuy nhiên trang web ngoài không cho biết số sinh viên đã đăng ký trong lớp, không cho biết điều kiện học phần và một số thông tin khác, khiến sinh viên cần mở nhiều tab một lúc cho nhiều ứng dụng để xem các thông tin khác nhau, từ đó xếp thời khóa biểu cho riêng mình.
\end{itemize}
Vì vậy, hệ thống sẽ tổng hợp các nhu cầu cá nhân cũng như thành tích học tập (mức cảnh cáo, các môn trượt) để đưa ra các gợi ý lựa chọn hợp lý nhất cho sinh viên.

\section{Tổng quan chức năng}
\subsection{Biểu đồ usecase tổng quát}
Biểu đồ usecase tổng quát mô tả quan hệ giữa các tác nhân với các chức năng chính của hệ thống. Các tác nhân tham gia vào hệ thống bao gồm: 

\begin{itemize}
    \item Admin: Quản lý tài khoản sinh viên, giáo viên, quản lý thông tin học phần, lớp mở.
    \item Sinh viên: Cập nhật thông tin các môn học đã học, trao đổi với chatbot về giải pháp đăng ký học tập.
\end{itemize}	

\begin{figure}[H]
    \centering
    \includegraphics[width=0.8\textwidth]{ảnh latex/usecase_tong_quat.png}
    \caption{Biểu đồ usecase tổng quát}
\end{figure}

\subsection{Biểu đồ usecase phân rã}
    \textbf{Biểu đồ 1: Theo dõi thông tin học tập}


Sinh viên theo dõi thông tin về điểm số các môn đã học, các môn đã trượt, các môn điểm kém cần cải thiện, các môn học nhanh hoặc chậm hơn so với chương trình đào tạo chuẩn, các học phần bắt buộc, các môn chưa học còn lại trong chương trình.

\begin{figure}[H]
    \centering
    \includegraphics[width=0.8\textwidth]{ảnh latex/usecase_theo_doi_ket_qua_hoc_tap.png}
    \caption{Biểu đồ usecase Theo dõi thông tin học tập}
\end{figure}

\textbf{Biểu đồ 2: Xây dựng thời khóa biểu mẫu}


Sinh viên có thể chọn các môn muốn học kỳ tới, hệ thống sẽ tự động trực quan thời khóa biểu đó cho sinh viên, đồng thời đảm bảo sinh viên không thể chọn các lớp trùng học phần hoặc trùng thời gian để tránh sai sót trong việc đăng ký lớp.
\begin{figure}[H]
    \centering
    \includegraphics[width=0.8\textwidth]{ảnh latex/usecase_xay_dung_thoi_khoa_bieu_mau.png}
    \caption{Biểu đồ usecase Xây dựng thời khóa biểu mẫu}
\end{figure}
\textbf{Biểu đồ 3: Trao đổi với chatbot}

Sinh viên có thể trao đổi với hệ thống chatbot để được đưa ra những tư vấn về các học phần và lớp học nên học dựa theo điểm số, các môn chưa học, tình trạng nợ môn và cảnh cáo học tập của sinh viên đó. [cite: 91]

\begin{figure}[H]
    \centering
    \includegraphics[width=0.8\textwidth]{ảnh latex/usecase_trao_doi_voi_chatbot.png}
    \caption{Biểu đồ phân rã usecase Trao đổi với chatbot}
\end{figure}

\textbf{Biểu đồ 4: Quản lý tài khoản sinh viên}

Admin có thể theo dõi, thêm, cập nhật, xóa tài khoản cho các sinh viên. [cite: 94]

\begin{figure}[H]
    \centering
    \includegraphics[width=0.8\textwidth]{ảnh latex/usecase_quan_ly_tai_khoan_sinh_vien.png}
    \caption{Biểu đồ phân rã usecase Quản lý tài khoản sinh viên}
\end{figure}

\textbf{Biểu đồ 5: Quản lý danh sách học phần}

Admin theo dõi, kiểm soát thông tin về các học phần trong chương trình đào tạo của từng ngành học. [cite: 97]

\begin{figure}[H]
    \centering
    \includegraphics[width=0.8\textwidth]{ảnh latex/usecase_quan_ly_danh_sanh_hoc_phan.png}
    \caption{Biểu đồ phân rã usecase Quản lý danh sách học phần}
\end{figure}

\textbf{Biểu đồ 6: Quản lý danh sách lớp mở}

Admin theo dõi, cập nhật thông tin các lớp sẽ được mở trong học kỳ mới. [cite: 100]

\begin{figure}[H]
    \centering
    \includegraphics[width=0.8\textwidth]{ảnh latex/usecase_quan_ly_danh_sach_lop_mo.png}
    \caption{Biểu đồ phân rã usecase Quản lý danh sách lớp mở}
\end{figure}

\textbf{Biểu đồ 7: Cập nhật điểm}

Sinh viên cập nhật danh sách điểm của mình lên hệ thống làm cơ sở để hệ thống đưa ra gợi ý đăng ký học tập cho sinh viên. [cite: 103]

\begin{figure}[H]
    \centering
    \includegraphics[width=0.8\textwidth]{ảnh latex/usecase_cap_nhat_diem.png}
    \caption{Biểu đồ phân rã usecase Cập nhật điểm}
\end{figure}

\subsection{Quy trình nghiệp vụ}
Quy trình cập nhật điểm số
\begin{figure}[H]
    \centering
    \includegraphics[width=0.9\textwidth]{ảnh latex/quy_trinh_cap_nhat_diem_so.png}
\end{figure}
Quy trình quản lý tài khoản sinh viên
\begin{figure}[H]
    \centering
    \includegraphics[width=0.9\textwidth]{ảnh latex/quy_trinh_quan_ly_tai_khoan_sinh_vien.png}
\end{figure}
Quy trình quản lý danh sách lớp mở
\begin{figure}[H]
    \centering
    \includegraphics[width=0.9\textwidth]{ảnh latex/quy_trinh_cap_nhat_lop.png}
\end{figure}


\subsection{Đặc tả chức năng}
\subsubsection{Đặc tả usecase: Xem thông tin học tập}

\begin{table}[H]
\renewcommand{\arraystretch}{1.4}
\centering
\begin{tabularx}{\textwidth}{|>{\columncolor[HTML]{C6E0B4}}l|l|l|X|}
\hline
\textbf{Mã Use case} & UC001 & \textbf{Tên Use case} & Xem thông tin học tập \\ \hline
\textbf{Tác nhân} & \multicolumn{3}{l|}{Sinh viên} \\ \hline
\textbf{Tiền điều kiện} & \multicolumn{3}{l|}{Sinh viên đã được cấp tài khoản và đăng nhập thành công vào hệ thống} \\ \hline
 & \multicolumn{3}{g|}{} \\ \pagecolor{white} % Reset màu cho dòng tiêu đề phụ
 & \cellcolor[HTML]{F8CBAD}\textbf{STT} & \cellcolor[HTML]{F8CBAD}\textbf{Thực hiện bởi} & \cellcolor[HTML]{F8CBAD}\textbf{Hành động} \\ \cline{2-4} 
\textbf{Luồng sự kiện} & 1. & Hệ thống & hiển thị màn hình chính gồm các mục: 
 \begin{itemize}[leftmargin=*, nosep]
     \item Xem kết quả học tập
     \item Xem chương trình đào tạo
 \end{itemize} \\ \cline{2-4} 
\textbf{chính} & 2. & Sinh viên & chọn một trong các mục cần xem \\ \cline{2-4} 
\textbf{(Thành công)} & 3. & Hệ thống & kiểm tra tính khả dụng của dữ liệu \\ \cline{2-4} 
 & 4. & Hệ thống & hiển thị thông tin tương ứng với danh mục đã chọn \\ \hline
 & \multicolumn{3}{g|}{} \\ 
\textbf{Luồng sự kiện} & \cellcolor[HTML]{F8CBAD}\textbf{STT} & \cellcolor[HTML]{F8CBAD}\textbf{Thực hiện bởi} & \cellcolor[HTML]{F8CBAD}\textbf{Hành động} \\ \cline{2-4} 
\textbf{thay thế} & 3a. & Hệ thống & nếu dữ liệu không khả dụng, thông báo: hệ thống đang bảo trì, vui lòng thử lại sau ít phút \\ \hline
\textbf{Hậu điều kiện} & \multicolumn{3}{l|}{Không có thay đổi dữ liệu sau khi xem thông tin} \\ \hline
\end{tabularx}
\caption{Đặc tả usecase Xem thông tin học tập}
\end{table}



\subsubsection{Đặc tả usecase: Xây dựng thời khóa biểu mẫu}

\begin{table}[H]
\renewcommand{\arraystretch}{1.4}
\centering
\begin{tabularx}{\textwidth}{|>{\columncolor[HTML]{C6E0B4}}l|l|l|X|}
\hline
\textbf{Mã Use case} & UC002 & \textbf{Tên Use case} & Xây dựng thời khóa biểu mẫu \\ \hline
\textbf{Tác nhân} & \multicolumn{3}{l|}{Sinh viên} \\ \hline
\textbf{Tiền điều kiện} & \multicolumn{3}{p{13cm}|}{Sinh viên đã được cấp tài khoản và đăng nhập thành công vào hệ thống} \\ \hline
 & \multicolumn{3}{l|}{} \\
 & \cellcolor[HTML]{F8CBAD}\textbf{STT} & \cellcolor[HTML]{F8CBAD}\textbf{Thực hiện bởi} & \cellcolor[HTML]{F8CBAD}\textbf{Hành động} \\ \cline{2-4} 
\textbf{Luồng sự kiện} & 1. & Hệ thống & hiển thị màn hình chính \\ \cline{2-4} 
\textbf{chính} & 2. & Sinh viên & chọn mục Xem thời khóa biểu \\ \cline{2-4} 
\textbf{(Thành công)} & 3. & Hệ thống & hiển thị màn hình Xem thời khóa biểu với danh sách lớp đã đăng ký và danh sách lớp khả dụng cho kỳ mới \\ \cline{2-4} 
 & 4. & Sinh viên & chọn lớp học muốn đăng ký và đăng ký hoặc chọn lớp đã có trong danh sách đã đăng ký để xóa \\ \cline{2-4} 
\multirow{-7}{*}{\textbf{}} & 5. & Hệ thống & cập nhật lớp vào thời khóa biểu theo tuần hiển thị ra cho sinh viên \\ \hline
 & \multicolumn{3}{l|}{} \\
\textbf{Luồng sự kiện} & \cellcolor[HTML]{F8CBAD}\textbf{STT} & \cellcolor[HTML]{F8CBAD}\textbf{Thực hiện bởi} & \cellcolor[HTML]{F8CBAD}\textbf{Hành động} \\ \cline{2-4} 
\textbf{thay thế} & 3a. & Hệ thống & nếu dữ liệu không khả dụng, thông báo: hệ thống đang bảo trì, vui lòng thử lại sau ít phút \\ \hline
\textbf{Hậu điều kiện} & \multicolumn{3}{l|}{Không có thay đổi dữ liệu sau khi xem thông tin} \\ \hline
\end{tabularx}
\caption{Đặc tả usecase Xây dựng thời khóa biểu mẫu}
\end{table}

\subsubsection{Đặc tả usecase: Hỏi chatbot tư vấn}

\begin{table}[H]
\renewcommand{\arraystretch}{1.4} % Giãn dòng để bảng dễ đọc
\centering
\begin{tabularx}{\textwidth}{|>{\columncolor[HTML]{C6E0B4}}l|l|l|X|}
\hline
\textbf{Mã Use case} & UC003 & \textbf{Tên Use case} & Hỏi chatbot tư vấn \\ \hline
\textbf{Tác nhân} & \multicolumn{3}{l|}{Sinh viên} \\ \hline
\textbf{Tiền điều kiện} & \multicolumn{3}{p{13cm}|}{Sinh viên đã được cấp tài khoản và đăng nhập thành công vào hệ thống} \\ \hline
 & \multicolumn{3}{l|}{} \\
 & \cellcolor[HTML]{F8CBAD}\textbf{STT} & \cellcolor[HTML]{F8CBAD}\textbf{Thực hiện bởi} & \cellcolor[HTML]{F8CBAD}\textbf{Hành động} \\ \cline{2-4} 
\textbf{Luồng sự kiện} & 1. & Hệ thống & hiển thị hộp thoại chatbot \\ \cline{2-4} 
\textbf{chính} & 2. & Sinh viên & ấn vào biểu tượng hộp thoại để chatbot hiện lên \\ \cline{2-4} 
\textbf{(Thành công)} & 3. & Sinh viên & đặt câu hỏi cho chatbot \\ \cline{2-4} 
\multirow{-6}{*}{\textbf{}} & 4. & Hệ thống & trả lời câu hỏi cho sinh viên \\ \hline
 & \multicolumn{3}{l|}{} \\
\textbf{Luồng sự kiện} & \cellcolor[HTML]{F8CBAD}\textbf{STT} & \cellcolor[HTML]{F8CBAD}\textbf{Thực hiện bởi} & \cellcolor[HTML]{F8CBAD}\textbf{Hành động} \\ \cline{2-4} 
\textbf{thay thế} & & & \\ \hline
\textbf{Hậu điều kiện} & \multicolumn{3}{l|}{Không có thay đổi dữ liệu sau khi xem thông tin} \\ \hline
\end{tabularx}
\caption{Đặc tả usecase Hỏi chatbot tư vấn}
\end{table}

\subsubsection{Đặc tả usecase: Thêm tài khoản sinh viên}

% Cấu hình màu sắc (Nếu đã khai báo ở đầu file thì không cần lặp lại)
\definecolor{lightgreen}{HTML}{C6E0B4}
\definecolor{lightorange}{HTML}{F8CBAD}

% --- BẢNG UC004: THÊM TÀI KHOẢN SINH VIÊN ---
\renewcommand{\arraystretch}{1.4}
\begin{xltabular}{\textwidth}{|>{\columncolor{lightgreen}\bfseries}l|l|l|X|}
\caption{Đặc tả usecase Thêm tài khoản sinh viên} \label{tab:uc004} \\
\hline
\rowcolor{lightgreen} Mã Use case & UC004 & Tên Use case & Thêm tài khoản sinh viên \\ \hline
Tác nhân & \multicolumn{3}{l|}{Admin} \\ \hline
Tiền điều kiện & \multicolumn{3}{l|}{Admin đăng nhập vào hệ thống} \\ \hline

\endhead 

% Luồng sự kiện chính
\rowcolor{lightgreen} \multicolumn{4}{|l|}{\textbf{Luồng sự kiện chính (Thành công)}} \\ \hline
\rowcolor{lightorange} & \textbf{STT} & \textbf{Thực hiện bởi} & \textbf{Hành động} \\ \hline
 & 1. & Hệ thống & hiển thị màn hình chính với nút Thêm tài khoản sinh viên \\ \cline{2-4}
 & 2. & Admin & ấn nút Thêm tài khoản sinh viên \\ \cline{2-4}
 & 3. & Hệ thống & hiển thị giao diện với các nút Upload file excel, Thêm thủ công và danh sách sinh viên đã tạo. \\ \cline{2-4}
 & 4a. & Admin & chọn Upload file excel \\ \cline{2-4}
 & 5a. & Hệ thống & hiển thị giao diện chọn file trong máy \\ \cline{2-4}
 & 6a. & Admin & chọn file Danh sách sinh viên, sau đó ấn Tạo \\ \cline{2-4}
 & 7a. & Hệ thống & cập nhật danh sách sinh viên và tự động cấp tài khoản cho các sinh viên trong danh sách \\ \cline{2-4}
 & 4b. & Admin & chọn Thêm thủ công \\ \cline{2-4}
 & 5b. & Hệ thống & hiển thị form điền thông tin sinh viên cần tạo \\ \cline{2-4}
 & 6b. & Admin & điền thông tin rồi ấn Tạo \\ \cline{2-4}
 & 7b. & Hệ thống & cập nhật danh sách sinh viên và tự động cấp tài khoản cho sinh viên mới tạo \\ \cline{2-4}
 & 4c. & Admin & chọn tài khoản cần sửa \\ \cline{2-4}
 & 5c. & Hệ thống & hiển thị thông tin tài khoản cần sửa \\ \cline{2-4}
 & 6c. & Admin & điền thông tin mới rồi ấn Sửa \\ \cline{2-4}
 & 7c. & Hệ thống & cập nhật danh sách sinh viên \\ \cline{2-4}
 & 4d. & Admin & chọn tài khoản cần xóa \\ \cline{2-4}
 & 5d. & Hệ thống & hiển thị hộp thoại xác nhận \\ \cline{2-4}
 & 6d. & Admin & xác nhận xóa tài khoản \\ \cline{2-4}
 & 7d. & Hệ thống & cập nhật danh sách sinh viên \\ \hline

% Luồng sự kiện thay thế
\rowcolor{lightgreen} \multicolumn{4}{|l|}{\textbf{Luồng sự kiện thay thế}} \\ \hline
\rowcolor{lightorange} & \textbf{STT} & \textbf{Thực hiện bởi} & \textbf{Hành động} \\ \hline
 & 6a.1 & Admin & nếu file không đúng định dạng, thông báo “File không đúng định dạng” và yêu cầu chọn lại \\ \cline{2-4}
 & 6a.2 & Admin & nếu Admin chọn Cancel, quay về 3 \\ \cline{2-4}
 & 6b.1 & Admin & nếu điền không đủ hoặc sai định dạng thông tin sinh viên, thông báo lỗi và yêu cầu điền đủ \\ \cline{2-4}
 & 6b.2 & Admin & nếu Admin chọn Cancel, quay về 3 \\ \cline{2-4}
 & 6c.1 & Admin & nếu điền không đủ hoặc sai định dạng thông tin sinh viên, thông báo lỗi và yêu cầu điền đủ \\ \cline{2-4}
 & 6c.2 & Admin & nếu Admin chọn Cancel, quay về 3 \\ \cline{2-4}
 & 6d.1 & Admin & nếu Admin chọn Cancel, quay về 3 \\ \hline

\rowcolor{lightgreen} Hậu điều kiện & \multicolumn{3}{l|}{Thêm danh sách sinh viên trên database và cấp tài khoản cho sinh viên mới thêm.} \\ \hline
\end{xltabular}

\subsubsection{Đặc tả usecase: Cập nhật điểm}

\definecolor{lightgreen}{HTML}{C6E0B4}
\definecolor{lightorange}{HTML}{F8CBAD}

\begin{xltabular}{\textwidth}{|>{\columncolor{lightgreen}\bfseries}l|l|l|X|}
\caption{Đặc tả usecase Cập nhật điểm} \label{tab:uc005} \\
\hline
\rowcolor{lightgreen} Mã Use case & UC005 & Tên Use case & Cập nhật điểm \\ \hline
Tác nhân & \multicolumn{3}{l|}{Sinh viên} \\ \hline
Tiền điều kiện & \multicolumn{3}{l|}{Sinh viên đăng nhập vào hệ thống} \\ \hline

\endhead 

% Luồng sự kiện chính
\rowcolor{lightgreen} \multicolumn{4}{|l|}{\textbf{Luồng sự kiện chính (Thành công)}} \\ \hline
\rowcolor{lightorange} & \textbf{STT} & \textbf{Thực hiện bởi} & \textbf{Hành động} \\ \hline
 & 1. & Hệ thống & hiển thị màn hình chính với nút Xem điểm \\ \cline{2-4}
 & 2. & Sinh viên & ấn nút Xem điểm \\ \cline{2-4}
 & 3. & Hệ thống & hiển thị giao diện với các nút Upload file excel, Thêm môn học \\ \cline{2-4}
 & 4a. & Sinh viên & chọn Upload file excel \\ \cline{2-4}
 & 5a. & Hệ thống & hiển thị giao diện chọn file trong máy \\ \cline{2-4}
 & 6a. & Sinh viên & chọn file điểm số, sau đó ấn Tạo \\ \cline{2-4}
 & 7a. & Hệ thống & cập nhật danh sách các môn đã học \\ \cline{2-4}
 & 4b. & Sinh viên & chọn Thêm môn học \\ \cline{2-4}
 & 5b. & Hệ thống & hiển thị form điền thông tin môn học đã học cần tạo \\ \cline{2-4}
 & 6b. & Sinh viên & điền thông tin rồi ấn Tạo \\ \cline{2-4}
 & 7b. & Hệ thống & cập nhật danh sách môn học \\ \cline{2-4}
 & 4c. & Sinh viên & chọn môn học cần xóa \\ \cline{2-4}
 & 5c. & Hệ thống & hiển thị thông báo xác nhận \\ \cline{2-4}
 & 6c. & Sinh viên & xác nhận xóa môn học \\ \cline{2-4}
 & 7c. & Hệ thống & cập nhật danh sách môn học đã học \\ \hline

% Luồng sự kiện thay thế
\rowcolor{lightgreen} \multicolumn{4}{|l|}{\textbf{Luồng sự kiện thay thế}} \\ \hline
\rowcolor{lightorange} & \textbf{STT} & \textbf{Thực hiện bởi} & \textbf{Hành động} \\ \hline
 & 6a.1 & Sinh viên & nếu file không đúng định dạng, thông báo “File không đúng định dạng” và yêu cầu chọn lại \\ \cline{2-4}
 & 6a.2 & Sinh viên & nếu sinh viên chọn Cancel, quay về 3 \\ \cline{2-4}
 & 6b.1 & Sinh viên & nếu điền không đủ hoặc sai định dạng, thông báo lỗi và yêu cầu điền đủ \\ \cline{2-4}
 & 6b.2 & Sinh viên & nếu Admin chọn Cancel, quay về 3 \\ \cline{2-4}
 & 6c.1 & Sinh viên & nếu điền không đủ hoặc sai định dạng, thông báo lỗi và yêu cầu điền đủ \\ \cline{2-4}
 & 6c.2 & Sinh viên & nếu sinh viên chọn Cancel, quay về 3 \\ \cline{2-4}
 & 6d.1 & Sinh viên & nếu sinh viên chọn Cancel, quay về 3 \\ \hline

\rowcolor{lightgreen} Hậu điều kiện & \multicolumn{3}{X|}{Cập nhật danh sách các môn đã học của sinh viên vào database và hiển thị lên giao diện.} \\ \hline
\end{xltabular}

\subsubsection{Đặc tả usecase: Cập nhật danh sách lớp mở}

\definecolor{lightgreen}{HTML}{C6E0B4}
\definecolor{lightorange}{HTML}{F8CBAD}

\begin{xltabular}{\textwidth}{|>{\columncolor{lightgreen}\bfseries}l|l|l|X|}
\caption{Đặc tả usecase Cập nhật danh sách lớp mở} \label{tab:uc006} \\
\hline
\rowcolor{lightgreen} Mã Use case & UC006 & Tên Use case & Cập nhật danh sách lớp mở \\ \hline
Tác nhân & \multicolumn{3}{l|}{Admin} \\ \hline
Tiền điều kiện & \multicolumn{3}{l|}{Admin đăng nhập vào hệ thống} \\ \hline

\endhead 

% Luồng sự kiện chính
\rowcolor{lightgreen} \multicolumn{4}{|l|}{\textbf{Luồng sự kiện chính (Thành công)}} \\ \hline
\rowcolor{lightorange} & \textbf{STT} & \textbf{Thực hiện bởi} & \textbf{Hành động} \\ \hline
 & 1. & Hệ thống & hiển thị màn hình chính với nút Cập nhật danh sách lớp mở \\ \cline{2-4}
 & 2. & Admin & ấn nút Cập nhật danh sách lớp mở \\ \cline{2-4}
 & 3. & Hệ thống & hiển thị giao diện với các nút Upload file excel, Thêm thủ công và danh sách lớp đã tạo. \\ \cline{2-4}
 & 4a. & Admin & chọn Upload file excel \\ \cline{2-4}
 & 5a. & Hệ thống & hiển thị giao diện chọn file trong máy \\ \cline{2-4}
 & 6a. & Admin & chọn file Danh sách lớp mở, sau đó ấn Tạo \\ \cline{2-4}
 & 7a. & Hệ thống & cập nhật danh sách lớp \\ \cline{2-4}
 & 4b. & Admin & chọn Thêm thủ công \\ \cline{2-4}
 & 5b. & Hệ thống & hiển thị form điền thông tin lớp \\ \cline{2-4}
 & 6b. & Admin & điền thông tin rồi ấn Tạo \\ \cline{2-4}
 & 7b. & Hệ thống & cập nhật danh sách lớp \\ \cline{2-4}
 & 4c. & Admin & chọn lớp cần sửa \\ \cline{2-4}
 & 5c. & Hệ thống & hiển thị thông tin lớp cần sửa \\ \cline{2-4}
 & 6c. & Admin & điền thông tin mới rồi ấn Sửa \\ \cline{2-4}
 & 7c. & Hệ thống & cập nhật danh sách lớp \\ \cline{2-4}
 & 4d. & Admin & chọn lớp cần xóa \\ \cline{2-4}
 & 5d. & Hệ thống & hiển thị hộp thoại xác nhận \\ \cline{2-4}
 & 6d. & Admin & xác nhận xóa lớp \\ \cline{2-4}
 & 7d. & Hệ thống & cập nhật danh sách lớp \\ \hline

% Luồng sự kiện thay thế
\rowcolor{lightgreen} \multicolumn{4}{|l|}{\textbf{Luồng sự kiện thay thế}} \\ \hline
\rowcolor{lightorange} & \textbf{STT} & \textbf{Thực hiện bởi} & \textbf{Hành động} \\ \hline
 & 6a.1 & Admin & nếu file không đúng định dạng, thông báo “File không đúng định dạng” và yêu cầu chọn lại \\ \cline{2-4}
 & 6a.2 & Admin & nếu Admin chọn Cancel, quay về 3 \\ \cline{2-4}
 & 6b.1 & Admin & nếu điền không đủ hoặc sai định dạng thông tin lớp, thông báo lỗi và yêu cầu điền đủ \\ \cline{2-4}
 & 6b.2 & Admin & nếu Admin chọn Cancel, quay về 3 \\ \cline{2-4}
 & 6c.1 & Admin & nếu điền không đủ hoặc sai định dạng thông tin lớp, thông báo lỗi và yêu cầu điền đủ \\ \cline{2-4}
 & 6c.2 & Admin & nếu Admin chọn Cancel, quay về 3 \\ \cline{2-4}
 & 6d.1 & Admin & nếu Admin chọn Cancel, quay về 3 \\ \hline

\rowcolor{lightgreen} Hậu điều kiện & \multicolumn{3}{X|}{Thêm danh sách lớp trên database.} \\ \hline
\end{xltabular}


\subsubsection{Đặc tả usecase: Cập nhật danh sách học phần}

\definecolor{lightgreen}{HTML}{C6E0B4}
\definecolor{lightorange}{HTML}{F8CBAD}

\begin{xltabular}{\textwidth}{|>{\columncolor{lightgreen}\bfseries}l|l|l|X|}
\caption{Đặc tả usecase Cập nhật danh sách học phần mở} \label{tab:uc006} \\
\hline
\rowcolor{lightgreen} Mã Use case & UC006 & Tên Use case & Cập nhật danh sách học phần \\ \hline
Tác nhân & \multicolumn{3}{l|}{Admin} \\ \hline
Tiền điều kiện & \multicolumn{3}{l|}{Admin đăng nhập vào hệ thống} \\ \hline

\endhead 

% Luồng sự kiện chính
\rowcolor{lightgreen} \multicolumn{4}{|l|}{\textbf{Luồng sự kiện chính (Thành công)}} \\ \hline
\rowcolor{lightorange} & \textbf{STT} & \textbf{Thực hiện bởi} & \textbf{Hành động} \\ \hline
 & 1. & Hệ thống & hiển thị màn hình chính với nút Cập nhật danh sách học phần \\ \cline{2-4}
 & 2. & Admin & ấn nút Cập nhật danh sách học phần \\ \cline{2-4}
 & 3. & Hệ thống & hiển thị giao diện với các nút Upload file excel, Thêm thủ công và danh sách học phần đã tạo. \\ \cline{2-4}
 & 4a. & Admin & chọn Upload file excel \\ \cline{2-4}
 & 5a. & Hệ thống & hiển thị giao diện chọn file trong máy \\ \cline{2-4}
 & 6a. & Admin & chọn file Danh sách học phần, sau đó ấn Tạo \\ \cline{2-4}
 & 7a. & Hệ thống & cập nhật danh sách học phần \\ \cline{2-4}
 & 4b. & Admin & chọn Thêm thủ công \\ \cline{2-4}
 & 5b. & Hệ thống & hiển thị form điền thông tin học phần \\ \cline{2-4}
 & 6b. & Admin & điền thông tin rồi ấn Tạo \\ \cline{2-4}
 & 7b. & Hệ thống & cập nhật danh sách học phần \\ \cline{2-4}
 & 4c. & Admin & chọn học phần cần sửa \\ \cline{2-4}
 & 5c. & Hệ thống & hiển thị thông tin học phần cần sửa \\ \cline{2-4}
 & 6c. & Admin & điền thông tin mới rồi ấn Sửa \\ \cline{2-4}
 & 7c. & Hệ thống & cập nhật danh sách học phần \\ \cline{2-4}
 & 4d. & Admin & chọn học phần cần xóa \\ \cline{2-4}
 & 5d. & Hệ thống & hiển thị hộp thoại xác nhận \\ \cline{2-4}
 & 6d. & Admin & xác nhận xóa học phần \\ \cline{2-4}
 & 7d. & Hệ thống & cập nhật danh sách học phần \\ \hline

% Luồng sự kiện thay thế
\rowcolor{lightgreen} \multicolumn{4}{|l|}{\textbf{Luồng sự kiện thay thế}} \\ \hline
\rowcolor{lightorange} & \textbf{STT} & \textbf{Thực hiện bởi} & \textbf{Hành động} \\ \hline
 & 6a.1 & Admin & nếu file không đúng định dạng, thông báo “File không đúng định dạng” và yêu cầu chọn lại \\ \cline{2-4}
 & 6a.2 & Admin & nếu Admin chọn Cancel, quay về 3 \\ \cline{2-4}
 & 6b.1 & Admin & nếu điền không đủ hoặc sai định dạng thông tin lớp, thông báo lỗi và yêu cầu điền đủ \\ \cline{2-4}
 & 6b.2 & Admin & nếu Admin chọn Cancel, quay về 3 \\ \cline{2-4}
 & 6c.1 & Admin & nếu điền không đủ hoặc sai định dạng thông tin lớp, thông báo lỗi và yêu cầu điền đủ \\ \cline{2-4}
 & 6c.2 & Admin & nếu Admin chọn Cancel, quay về 3 \\ \cline{2-4}
 & 6d.1 & Admin & nếu Admin chọn Cancel, quay về 3 \\ \hline

\rowcolor{lightgreen} Hậu điều kiện & \multicolumn{3}{X|}{Thêm danh sách học phần trên database.} \\ \hline
\end{xltabular}



\section{Yêu cầu phi chức năng}
% \section{Yêu cầu phi chức năng}

\subsection{Hiệu năng}
\begin{itemize}
    \item \textbf{Thời gian phản hồi:} Giao diện người dùng cần phản hồi trong vòng 2 giây với mỗi thao tác, các thao tác cần xử lý file (danh sách sinh viên, danh sách học phần, danh sách lớp học) cần phản hồi nhanh chóng tùy vào kích thước file.
    \item \textbf{Xử lý đồng thời:} Hệ thống có thể hỗ trợ tối thiểu 10000 yêu cầu đồng thời mà không ảnh hưởng đến hiệu suất.
\end{itemize}

\subsection{Tính sẵn sàng và phục hồi}
\begin{itemize}
    \item \textbf{Tính sẵn sàng:} Hệ thống hoạt động 24/7, thời gian ngừng hoạt động không quá 1 giờ mỗi tháng.
    \item \textbf{Khả năng phục hồi:} Có thể sao lưu dữ liệu hàng ngày và khôi phục khi có sự cố. Dữ liệu được lưu trữ tối thiểu 5 năm.
\end{itemize}

\subsection{Bảo mật}
\begin{itemize}
    \item \textbf{Xác thực và phân quyền:} Hệ thống yêu cầu đăng nhập bằng tài khoản riêng, phân quyền theo vai trò.
    \item \textbf{Mã hóa dữ liệu:} Thông tin nhạy cảm được mã hóa khi lưu trữ và truyền tải.
    \item \textbf{Ghi nhật ký hệ thống:} Ghi lại các hành động quan trọng như đăng nhập, cập nhật, truy xuất dữ liệu.
\end{itemize}

\subsection{Khả năng mở rộng}
\begin{itemize}
    \item \textbf{Kiến trúc linh hoạt:} Cho phép mở rộng dễ dàng theo số lượng người dùng và các module chức năng mới.
\end{itemize}

\subsection{Tính khả dụng}
\begin{itemize}
    \item \textbf{Giao diện thân thiện:} Giao diện dễ sử dụng cho cả sinh viên và admin.
    \item \textbf{Hỗ trợ nền tảng web:} Hệ thống hoạt động tốt trên trình duyệt web.
\end{itemize}

\subsection{Độ tin cậy}
\begin{itemize}
    \item \textbf{Toàn vẹn dữ liệu:} Hệ thống đảm bảo dữ liệu được xử lý chính xác, không trùng lặp, không mất mát.
    \item \textbf{Kiểm thử toàn diện:} Trước khi triển khai, hệ thống phải được kiểm thử nghiêm ngặt để đảm bảo độ tin cậy.
\end{itemize}

\subsection{Tương thích}
\begin{itemize}
    \item \textbf{Trình duyệt phổ biến:} Hệ thống hoạt động tốt trên Chrome, Firefox, Safari và Edge.
    \item \textbf{Thiết bị di động:} Hỗ trợ giao diện responsive trên Android và iOS.
\end{itemize}

\subsection{Khả năng bảo trì}
\begin{itemize}
    \item \textbf{Kiến trúc rõ ràng:} Mã nguồn được tổ chức dễ hiểu, dễ nâng cấp và bảo trì.
    \item \textbf{Cập nhật linh hoạt:} Hệ thống có thể cập nhật chức năng hoặc sửa lỗi mà không làm gián đoạn dịch vụ.
\end{itemize}

% --- CHƯƠNG 3 ---
\chapter{Công nghệ sử dụng}

\section{Kiến trúc hệ thống}
Dự án được triển khai theo mô hình \textbf{3 tầng} gồm giao diện, xử lý nghiệp vụ và cơ sở dữ liệu. Cách tổ chức này giúp hệ thống dễ mở rộng, dễ kiểm thử và tách biệt rõ ràng trách nhiệm của từng lớp.

\begin{itemize}
    \item \textbf{Tầng giao diện:} xây dựng bằng React 19 và TypeScript, chịu trách nhiệm hiển thị màn hình cho hai nhóm người dùng chính là sinh viên và quản trị viên.
    \item \textbf{Tầng dịch vụ:} FastAPI đảm nhiệm việc cung cấp REST API, xác thực người dùng, xử lý nghiệp vụ và điều phối các dịch vụ hỗ trợ như chatbot, xuất Excel, gửi email và tính toán điểm.
    \item \textbf{Tầng dữ liệu:} MySQL lưu trữ dữ liệu học tập, tài khoản, thời khóa biểu, lịch sử đăng ký lớp, lịch sử hội thoại chatbot và các biểu mẫu hỗ trợ.
\end{itemize}

Ngoài luồng xử lý đồng bộ truyền thống, hệ thống còn có các cơ chế bổ trợ như lưu lịch sử hội thoại, streaming phản hồi cho chatbot và hàng đợi xử lý nội bộ cho một số tác vụ nền. Nhờ đó, ứng dụng có thể phản hồi nhanh nhưng vẫn giữ được khả năng mở rộng trong tương lai.

\section{Công nghệ backend}
\subsection{FastAPI}
FastAPI được sử dụng làm khung backend chính vì hỗ trợ tốt kiểu lập trình bất đồng bộ, có hiệu năng cao và tự sinh tài liệu API.
\begin{itemize}
    \item Hệ thống endpoint được tổ chức rõ ràng theo từng nhóm chức năng như sinh viên, học phần, lớp học, điểm, xác thực, chatbot và quản trị.
    \item Swagger UI và ReDoc giúp việc kiểm thử API, tích hợp frontend và rà soát lỗi nghiệp vụ thuận tiện hơn.
\end{itemize}

\subsection{SQLAlchemy và MySQL}
\begin{itemize}
    \item \textbf{SQLAlchemy 2.0.25} được dùng để ánh xạ dữ liệu theo mô hình đối tượng, hạn chế viết SQL thuần và giảm nguy cơ lỗi truy vấn.
    \item \textbf{PyMySQL 1.0.2} là driver kết nối với MySQL 8.0.
    \item Cấu hình kết nối được đặt tập trung trong \texttt{backend/app/core/config.py} và \texttt{backend/app/db/database.py}, giúp quản lý môi trường dev và production dễ hơn.
\end{itemize}

\subsection{Xác thực và bảo mật}
\begin{itemize}
    \item \textbf{JWT}: dùng để phát hành và kiểm tra token đăng nhập.
    \item \textbf{Argon2 / passlib}: băm mật khẩu trước khi lưu vào CSDL.
    \item \textbf{CORS}: cấu hình theo danh sách domain được phép truy cập.
    \item \textbf{Rate limiting}: hạn chế một số tác vụ nhạy cảm như khôi phục mật khẩu để giảm rủi ro spam và brute-force.
\end{itemize}

\subsection{Dịch vụ hỗ trợ}
\begin{itemize}
    \item \textbf{Gửi email}: phục vụ xác thực và khôi phục mật khẩu.
    \item \textbf{Xuất Excel}: hỗ trợ tải xuống các phương án thời khóa biểu và các bảng dữ liệu phục vụ người dùng.
    \item \textbf{Tính toán điểm}: tính CPA/GPA, tổng tín chỉ và mức cảnh báo học tập từ dữ liệu học phần đã học.
\end{itemize}

\section{Công nghệ chatbot}
\subsection{Phân loại ý định}
Chatbot sử dụng bộ phân loại lai để nhận diện ý định người dùng dựa trên dữ liệu tiếng Việt.
\begin{itemize}
    \item \textbf{TF-IDF + cosine similarity}: đo độ gần nhau giữa câu hỏi hiện tại và các mẫu huấn luyện.
    \item \textbf{Word2Vec}: bổ sung góc nhìn ngữ nghĩa, giúp hiểu những cách diễn đạt khác nhau nhưng cùng mục đích.
    \item \textbf{Điểm từ khóa và điểm khớp mẫu}: tăng độ chính xác cho các câu hỏi ngắn hoặc có mẫu rõ ràng.
    \item \textbf{Adaptive weights}: thay đổi trọng số theo độ dài câu hỏi để phù hợp với từng kiểu truy vấn.
\end{itemize}

\subsection{Luồng xử lý hội thoại}
Luồng xử lý chatbot trong backend gồm các bước chính:
\begin{itemize}
    \item Chuẩn hóa văn bản đầu vào.
    \item Phân loại ý định và tách câu hỏi nhiều phần nếu cần.
    \item Rẽ nhánh sang xử lý rule-based, NL2SQL hoặc agent orchestration tùy theo intent và cấu hình môi trường.
    \item Trả về phản hồi dạng văn bản, bảng dữ liệu hoặc tổ hợp lớp học kèm metadata để frontend hiển thị tương tác.
\end{itemize}

\subsection{Rule engine cho gợi ý học phần và lớp học}
Hai bộ luật chính được triển khai trong \texttt{SubjectSuggestionRuleEngine} và \texttt{ClassSuggestionRuleEngine}.
\begin{itemize}
    \item \textbf{Gợi ý học phần:} ưu tiên các môn cần học lại, môn đúng lộ trình chương trình, môn phù hợp tín chỉ còn lại và tình trạng học tập của sinh viên.
    \item \textbf{Gợi ý lớp học:} lọc theo xung đột thời gian, số lượng lớp trùng môn, ngày học mong muốn, khung giờ sáng/chiều, giáo viên ưu tiên và các ràng buộc mềm khác.
    \item \textbf{Tạo tổ hợp lớp:} dùng phép tổ hợp để tạo nhiều phương án thời khóa biểu khả dĩ rồi chấm điểm theo mức độ phù hợp với sở thích người dùng.
\end{itemize}

\subsection{NL2SQL và tra cứu dữ liệu}
\begin{itemize}
    \item NL2SQL chuyển những câu hỏi như xem điểm, xem thông tin học phần, xem lớp học hoặc xem lịch học thành truy vấn dữ liệu có cấu trúc.
    \item Câu hỏi được trích xuất thực thể bằng biểu thức chính quy, sau đó ghép với mẫu truy vấn phù hợp trong kho ví dụ có sẵn.
    \item Với các truy vấn lớp học, hệ thống còn dùng bộ tách ràng buộc và truy vấn ORM an toàn để tránh lỗ hổng chèn lệnh SQL.
\end{itemize}

\subsection{Lưu lịch sử và streaming}
\begin{itemize}
    \item Chatbot lưu lịch sử hội thoại theo từng sinh viên, cho phép đổi tên, xóa và tải lại các cuộc trò chuyện trước đó.
    \item Phản hồi được stream theo các trạng thái như chuẩn hóa, phân loại, truy vấn và định dạng để giao diện hiển thị tiến trình rõ ràng hơn.
\end{itemize}

\section{Công nghệ frontend}
\subsection{Framework và công cụ}
\begin{itemize}
    \item \textbf{React 19.1.0} và \textbf{TypeScript 5.8.3} là nền tảng chính của frontend.
    \item \textbf{Vite 7.0.4} giúp khởi động và build nhanh.
    \item \textbf{React Router v7} quản lý điều hướng giữa các màn hình cho sinh viên và admin.
\end{itemize}

\subsection{Thư viện giao diện}
\begin{itemize}
    \item \textbf{Ant Design 5.27.3}: cung cấp layout, form, bảng, modal, popup và nhiều component quản trị phổ biến.
    \item \textbf{Tailwind CSS 4.1.11}: hỗ trợ bố cục responsive và tinh chỉnh giao diện nhanh.
    \item \textbf{Chart.js 4.5.0}: dùng để trực quan hóa dữ liệu học tập, điểm số và thống kê.
    \item \textbf{xlsx, docx, file-saver}: phục vụ xuất và tải xuống dữ liệu.
\end{itemize}

\subsection{Các màn hình chính}
Frontend được chia thành các nhóm chức năng rõ ràng:
\begin{itemize}
    \item \textbf{Sinh viên:} trang tổng quan, lịch học, điểm số, chương trình đào tạo, đăng ký học phần, đăng ký lớp, biểu mẫu, FAQ, phản hồi, đổi mật khẩu và hướng dẫn sử dụng.
    \item \textbf{Quản trị viên:} dashboard quản trị, quản lý sinh viên, học phần, lớp học, thời khóa biểu, phản hồi, FAQ, thông báo và đổi mật khẩu.
    \item \textbf{Chatbot UI:} giao diện hội thoại riêng có lưu lịch sử và hỗ trợ các lựa chọn tương tác khi chatbot cần thu thập thêm thông tin.
\end{itemize}

\section{Cơ sở dữ liệu}
Hệ thống sử dụng \textbf{MySQL 8.0}. Từ metadata của mô hình dữ liệu, dự án hiện có \textbf{19 bảng} chính, bao phủ các nhóm dữ liệu sau:
\begin{itemize}
    \item \textbf{Người dùng và xác thực:} \texttt{students}, \texttt{admins}, \texttt{password\_reset\_tokens}.
    \item \textbf{Tổ chức đào tạo:} \texttt{departments}, \texttt{courses}, \texttt{subjects}, \texttt{course\_subjects}, \texttt{student\_course}, \texttt{subject\_condition}.
    \item \textbf{Lớp học và đăng ký:} \texttt{classes}, \texttt{class\_registers}, \texttt{subject\_registers}, \texttt{linked\_class}.
    \item \textbf{Kết quả học tập:} \texttt{learned\_subjects}, \texttt{semester\_gpa}.
    \item \textbf{Hỗ trợ và hội thoại:} \texttt{feedbacks}, \texttt{faqs}, \texttt{chat\_conversations}, \texttt{chat\_messages}.
\end{itemize}

Thiết kế này cho phép hệ thống lưu đầy đủ thông tin về sinh viên, chương trình đào tạo, dữ liệu đăng ký, lịch sử học tập và lịch sử chatbot trong cùng một CSDL trung tâm.

\section{Công cụ phát triển và kiểm thử}
\begin{itemize}
    \item \textbf{Python 3.11}: ngôn ngữ chính của backend.
    \item \textbf{Node.js, npm}: môi trường dựng frontend.
    \item \textbf{Pytest}: kiểm thử các bộ luật, bộ phân loại và thành phần xử lý đăng ký lớp.
    \item \textbf{Uvicorn}: chạy FastAPI trong quá trình phát triển và triển khai.
    \item \textbf{ESLint / TypeScript compiler}: kiểm tra chất lượng mã nguồn frontend.
\end{itemize}

% --- CHƯƠNG 4 ---
\chapter{Thiết kế và triển khai hệ thống}

\section{Phát triển cơ sở dữ liệu}
Thiết kế cơ sở dữ liệu được xây dựng theo hướng chuẩn hóa dữ liệu, ưu tiên các bảng lõi phục vụ nghiệp vụ học tập và các bảng phụ trợ phục vụ chatbot, xác thực và lịch sử hội thoại.

\begin{itemize}
    \item Các bảng \texttt{students}, \texttt{courses}, \texttt{subjects} và \texttt{classes} là trục dữ liệu chính của hệ thống.
    \item Các bảng trung gian như \texttt{course\_subjects}, \texttt{class\_registers}, \texttt{subject\_registers} và \texttt{student\_course} cho phép biểu diễn quan hệ nhiều-nhiều giữa sinh viên, học phần và lớp học.
    \item Các bảng \texttt{learned\_subjects} và \texttt{semester\_gpa} lưu lịch sử học tập để tính điểm trung bình, tổng tín chỉ và mức cảnh báo học vụ.
    \item Các bảng \texttt{chat\_conversations} và \texttt{chat\_messages} lưu toàn bộ lịch sử trao đổi với chatbot, hỗ trợ tra cứu lại phiên làm việc cũ.
\end{itemize}

Việc tách bảng theo đúng vai trò giúp truy vấn dễ hơn, dữ liệu rõ nghĩa hơn và giảm rủi ro trùng lặp thông tin. Các ràng buộc khóa chính, khóa ngoại và quan hệ giữa bảng được khai báo ngay trong mô hình SQLAlchemy, nhờ đó việc tạo CSDL và đồng bộ schema trở nên thống nhất với code.

\section{Phát triển backend}
Backend được chia thành các lớp rõ ràng để dễ bảo trì:
\begin{itemize}
    \item \textbf{models}: định nghĩa cấu trúc bảng bằng SQLAlchemy.
    \item \textbf{schemas}: mô tả dữ liệu vào/ra bằng Pydantic.
    \item \textbf{routes}: khai báo API endpoint.
    \item \textbf{services}: chứa phần nghiệp vụ như chatbot, tính điểm, xuất Excel, tìm kiếm mờ, gửi email và quản lý lịch sử hội thoại.
    \item \textbf{rules}: lưu các bộ luật gợi ý học phần và lớp học.
\end{itemize}

Một số nhóm API và dịch vụ quan trọng:
\begin{itemize}
    \item \textbf{Xác thực}: đăng nhập, đăng ký, quên mật khẩu, đặt lại mật khẩu.
    \item \textbf{Sinh viên}: hồ sơ, điểm số, lịch học, chương trình đào tạo, thay đổi mật khẩu.
    \item \textbf{Học phần và lớp học}: tra cứu, thêm, sửa, xóa, đăng ký và hủy đăng ký.
    \item \textbf{Quản trị}: quản lý sinh viên, học phần, lớp, thông báo, phản hồi và FAQ.
    \item \textbf{Chatbot}: hội thoại, phân loại intent, gợi ý học phần, gợi ý lớp, xử lý truy vấn theo ngôn ngữ tự nhiên, lưu lịch sử chat và xuất phương án lịch học.
\end{itemize}

Ngoài các route chính, backend còn có lớp kiểm tra tình trạng hệ thống, hỗ trợ healthcheck, cache Redis và cơ chế điều phối agent tùy chọn qua biến môi trường.

\subsection{Thiết kế lớp}
Phần thiết kế lớp của hệ thống tập trung vào các thực thể lõi phản ánh trực tiếp nghiệp vụ đăng ký học tập của sinh viên. Trong phạm vi này, bốn lớp quan trọng nhất là \texttt{Student}, \texttt{Subject}, \texttt{Class} và \texttt{LearnedSubject}. Các lớp này được chọn vì chúng xuất hiện xuyên suốt cả ba luồng nghiệp vụ trọng tâm của hệ thống: xem kết quả học tập và upload bảng điểm, tư vấn học phần nên đăng ký, và tư vấn lớp học nên đăng ký.

Trong cách tổ chức hiện tại, phần lớn lớp miền được hiện thực bằng SQLAlchemy. Vì vậy, thay vì chỉ mô tả thuộc tính dữ liệu, báo cáo cũng làm rõ vai trò nghiệp vụ, các validator quan trọng và các hành vi thường xuyên đi kèm trong service/route. Cách trình bày này giúp thấy rõ lớp nào giữ dữ liệu, lớp nào xử lý quy tắc và lớp nào đóng vai trò điều phối giao tiếp giữa frontend với backend.

\begin{table}[H]
\centering
\renewcommand{\arraystretch}{1.35}
\begin{tabularx}{\textwidth}{|l|X|X|X|}
\hline
\rowcolor[HTML]{D9EAD3}
\textbf{Lớp} & \textbf{Thuộc tính chính} & \textbf{Phương thức / validator tiêu biểu} & \textbf{Vai trò trong hệ thống} \\ \hline
\texttt{Student} & \texttt{id}, \texttt{student\_name}, \texttt{email}, \texttt{course\_id}, \texttt{department\_id}, \texttt{cpa}, \texttt{warning\_level}, \texttt{year\_level}, \texttt{total\_learned\_credits}, \texttt{total\_failed\_credits} & Được truy xuất qua các route hồ sơ học tập; dữ liệu CPA, cảnh báo và năm học được cập nhật khi thay đổi \texttt{LearnedSubject} & Đại diện cho sinh viên, làm gốc cho các quan hệ đăng ký học phần, đăng ký lớp và thống kê học tập \\ \hline
\texttt{Subject} & \texttt{id}, \texttt{subject\_id}, \texttt{subject\_name}, \texttt{credits}, \texttt{duration}, \texttt{department\_id}, \texttt{conditional\_subjects} & Tra cứu theo mã học phần, liên kết với \texttt{CourseSubject}, \texttt{LearnedSubject}, \texttt{Class} và \texttt{SubjectRegister} & Đại diện cho học phần, là đầu vào của cả gợi ý học phần lẫn gợi ý lớp học \\ \hline
\texttt{Class} & \texttt{id}, \texttt{subject\_id}, \texttt{class\_id}, \texttt{class\_name}, \texttt{linked\_class\_ids}, \texttt{class\_type}, \texttt{classroom}, \texttt{study\_date}, \texttt{study\_time\_start}, \texttt{study\_time\_end}, \texttt{teacher\_name}, \texttt{study\_week} & \texttt{validate\_study\_date()}, \texttt{validate\_time()}, logic ghép các lớp có cùng mã lớp hoặc lớp kèm & Lưu thông tin lớp mở, phục vụ đăng ký lớp và sinh các phương án thời khóa biểu không trùng lịch \\ \hline
\texttt{LearnedSubject} & \texttt{id}, \texttt{subject\_name}, \texttt{credits}, \texttt{letter\_grade}, \texttt{semester}, \texttt{student\_id}, \texttt{subject\_id} & Được tạo/sửa/xóa qua các API học tập; kéo theo \texttt{update\_semester\_gpa()} và \texttt{update\_student\_stats()} & Lưu kết quả học tập theo học kỳ, là nguồn dữ liệu chính để tính GPA/CPA, cảnh báo học tập và gợi ý học phần cần đăng ký lại \\ \hline
\end{tabularx}
\caption{Các lớp chủ đạo trong phần thiết kế lớp}
\end{table}

\subsubsection{1. Lớp \texttt{Student}}
\texttt{Student} là lớp trung tâm của toàn bộ hệ thống. Một đối tượng sinh viên không chỉ chứa thông tin định danh như họ tên, email hay mã khoa, mà còn lưu các chỉ số học tập tổng hợp như CPA, số tín chỉ đã tích lũy, số tín chỉ trượt, mức cảnh báo và trình độ năm học. Các thuộc tính này xuất hiện trực tiếp trong màn hình dashboard của sinh viên và cũng là đầu vào cho chatbot khi gợi ý học phần hoặc lớp học.

Về quan hệ, \texttt{Student} liên kết một-nhiều với \texttt{LearnedSubject}, \texttt{SemesterGPA}, \texttt{ClassRegister} và \texttt{SubjectRegister}. Nhờ các liên kết này, hệ thống có thể nhanh chóng truy xuất lịch sử học tập, lịch sử đăng ký và các lớp đang theo học của sinh viên. Trong luồng \texttt{/students/\{student\_id\}/academic-details}, backend tổng hợp dữ liệu từ các quan hệ này để trả về một cấu trúc học tập đầy đủ cho frontend.

Ở góc độ nghiệp vụ, các hành vi cập nhật của \texttt{Student} chủ yếu đến từ service thay vì method tự thân của model. Khi sinh viên thêm, sửa hoặc xóa một môn trong \texttt{LearnedSubject}, các hàm như \texttt{update\_student\_stats()} sẽ tự động tính lại CPA, số môn trượt, số tín chỉ tích lũy và mức cảnh báo. Điều này giúp lớp \texttt{Student} đóng vai trò “hồ sơ học tập động” thay vì chỉ là bản ghi tĩnh trong cơ sở dữ liệu.

\subsubsection{2. Lớp \texttt{Subject}}
\texttt{Subject} biểu diễn một học phần trong chương trình đào tạo. Lớp này lưu mã học phần, tên học phần, số tín chỉ, thời lượng, khoa phụ trách và các điều kiện tiên quyết. Đây là lớp được dùng nhiều nhất khi hệ thống cần tra cứu học phần theo mã, đối chiếu với chương trình đào tạo của sinh viên hoặc sinh ra danh sách học phần nên đăng ký.

Trong luồng đăng ký học phần, \texttt{Subject} là điểm nối giữa chương trình đào tạo và dữ liệu học tập thực tế. Bảng \texttt{course\_subjects} dùng để xác định những học phần nào thuộc từng khóa học, còn \texttt{learned\_subjects} cho biết sinh viên đã học hay chưa học môn đó. Khi chatbot tư vấn, service sẽ kết hợp \texttt{Subject} với \texttt{LearnedSubject} để nhận biết môn nào chưa học, môn nào đã học trượt và môn nào là ứng viên phù hợp cho kỳ tới.

\texttt{Subject} cũng liên kết trực tiếp với \texttt{Class}. Điều này cho phép hệ thống đi từ một học phần sang danh sách lớp mở tương ứng, rồi từ đó lọc tiếp theo thời gian, giảng viên và phòng học. Nhờ vậy, lớp \texttt{Subject} giữ vai trò “điểm gốc” cho cả hai nhánh tư vấn: chọn học phần và chọn lớp.

\subsubsection{3. Lớp \texttt{Class}}
\texttt{Class} là lớp mô tả một lớp mở cụ thể của một học phần. Khác với \texttt{Subject} là khái niệm học phần trừu tượng, \texttt{Class} chứa thông tin vận hành thực tế như mã lớp, tên lớp, loại lớp, phòng học, lịch học, giảng viên, tuần học và danh sách lớp kèm. Đây là lớp quyết định trực tiếp trải nghiệm đăng ký lớp của sinh viên vì mọi phương án đăng ký đều phải thỏa mãn ràng buộc thời gian từ lớp này.

Hai validator quan trọng của \texttt{Class} là \texttt{validate\_study\_date()} và \texttt{validate\_time()}. Validator đầu tiên chuẩn hóa và kiểm tra chuỗi ngày học, bảo đảm chỉ nhận các giá trị hợp lệ như Monday, Tuesday, ...; validator thứ hai chuyển thời gian từ chuỗi sang kiểu \texttt{time} và loại bỏ các giá trị không hợp lệ. Cách làm này giúp dữ liệu lịch học luôn nhất quán ngay từ tầng model, tránh lỗi khi sinh thời khóa biểu hoặc so sánh xung đột giờ học.

Về quan hệ, \texttt{Class} liên kết nhiều-nhất với \texttt{ClassRegister} và nhiều-nhất với \texttt{Subject}. Khi sinh viên đăng ký lớp, frontend sẽ hiển thị danh sách lớp từ endpoint \texttt{/api/classes/}, còn backend sẽ ghi nhận bản ghi đăng ký tương ứng trong \texttt{class\_registers}. Trong phần gợi ý lớp học, service còn dùng thêm thuộc tính \texttt{linked\_class\_ids} để nhóm các buổi học cùng một lớp logic, bảo đảm một phương án thời khóa biểu không bị tách rời sai cách.

\subsubsection{4. Lớp \texttt{LearnedSubject}}
\texttt{LearnedSubject} lưu lại một kết quả học tập của sinh viên theo từng học kỳ. Mỗi bản ghi chứa tên học phần, số tín chỉ, điểm chữ, học kỳ, khóa ngoại tới sinh viên và khóa ngoại tới học phần. Đây là lớp có tác động lớn nhất đến các chỉ số học tập tổng hợp, vì nó là đầu vào cho tính GPA học kỳ, CPA toàn khóa và số tín chỉ trượt.

Luồng upload bảng điểm của hệ thống đi qua \texttt{LearnedSubject} trước tiên. Frontend ở trang \texttt{Grades.tsx} và component \texttt{GradeExcelUpload.tsx} đọc file Excel, trích từng dòng hợp lệ rồi gọi API tạo mới môn đã học. Ở backend, route \texttt{create\_new\_learned\_subject()} tìm \texttt{Subject} theo mã học phần, tạo \texttt{LearnedSubject}, sau đó gọi \texttt{update\_semester\_gpa()} và \texttt{update\_student\_stats()} để tính lại toàn bộ chỉ số liên quan. Vì vậy, \texttt{LearnedSubject} vừa là lớp lưu trữ lịch sử học tập, vừa là nguồn kích hoạt cho các phép tính thống kê tự động.

Khi một bản ghi \texttt{LearnedSubject} được thêm, sửa hoặc xóa, backend sẽ đồng thời cập nhật bảng \texttt{semester\_gpa} và các trường tổng hợp của \texttt{Student}. Điều này giúp dữ liệu điểm số và dữ liệu hồ sơ sinh viên luôn đồng bộ, đồng thời làm cơ sở cho chatbot gợi ý học phần nên học lại hoặc nên đăng ký tiếp ở kỳ kế tiếp.

Ngoài bốn lớp lõi trên, hệ thống còn có một số lớp dịch vụ trung tâm đóng vai trò điều phối nghiệp vụ. Trong đó, \texttt{ChatbotService} chịu trách nhiệm phân loại intent, gom dữ liệu từ \texttt{Student}, \texttt{Subject}, \texttt{LearnedSubject}, \texttt{Class} và các rule engine để tạo phản hồi cho frontend; \texttt{ClassQueryService} và \texttt{ScheduleCombinationService} lọc và ghép các lớp theo ràng buộc thời gian; \texttt{ExcelExportService} xuất phương án lịch học ra file Excel để sinh viên tải xuống. Các lớp này không được đưa vào bảng chi tiết vì chúng thiên về xử lý nghiệp vụ hơn là mô hình dữ liệu, nhưng lại xuất hiện trực tiếp trong ba luồng trình tự ở phần sau.

\begin{figure}[H]
\centering
\resizebox{\textwidth}{!}{%
\begin{tikzpicture}[>=Latex, font=\small]
    \node[draw, rounded corners, minimum width=2.4cm, minimum height=0.8cm, align=center] (sv) at (0,0) {Sinh viên};
    \node[draw, rounded corners, minimum width=3.4cm, minimum height=0.8cm, align=center] (ui) at (3.5,0) {GradeExcelUpload\\và trang Xem điểm};
    \node[draw, rounded corners, minimum width=4.2cm, minimum height=0.8cm, align=center] (api) at (7.8,0) {\texttt{/learned-subjects/create-new-learned-subject}};
    \node[draw, rounded corners, minimum width=2.8cm, minimum height=0.8cm, align=center] (subject) at (12.2,0) {\texttt{Subject}};
    \node[draw, rounded corners, minimum width=3.2cm, minimum height=0.8cm, align=center] (ls) at (16.1,0) {\texttt{LearnedSubject}\\và GPA};

    \draw[dashed] (0,-0.6) -- (0,-6.7);
    \draw[dashed] (3.5,-0.6) -- (3.5,-6.7);
    \draw[dashed] (7.8,-0.6) -- (7.8,-6.7);
    \draw[dashed] (12.2,-0.6) -- (12.2,-6.7);
    \draw[dashed] (16.1,-0.6) -- (16.1,-6.7);

    \draw[->] (0.7,-1.1) -- node[above, align=center]{Chọn file Excel\\và xác nhận upload} (3.1,-1.1);
    \draw[->] (3.1,-2.0) -- node[above, align=center]{Parse từng dòng hợp lệ\\rồi tạo request} (7.4,-2.0);
    \draw[->] (7.4,-3.0) -- node[above, align=center]{Tìm học phần theo\\mã học phần} (11.8,-3.0);
    \draw[->] (11.8,-4.0) -- node[above, align=center]{Tạo bản ghi\\\texttt{LearnedSubject}} (15.8,-4.0);
    \draw[->] (15.8,-5.0) -- node[above, align=center]{Cập nhật GPA học kỳ\\và thống kê sinh viên} (7.9,-5.0);
    \draw[->] (7.9,-5.9) -- node[above, align=center]{Trả kết quả thành công\\cho giao diện} (3.6,-5.9);
\end{tikzpicture}
}%
\caption{Luồng truyền thông điệp khi sinh viên upload bảng điểm}
\end{figure}

Trong luồng này, giao diện của sinh viên làm hai việc chính: đọc file Excel và gọi API tạo từng bản ghi \texttt{LearnedSubject}. Backend sau đó tự động tìm \texttt{Subject}, ghi dữ liệu và tính lại các chỉ số học tập. Cách thiết kế này phù hợp với logic FE hiện có, nơi component upload chỉ đóng vai trò nhập dữ liệu, còn toàn bộ kiểm tra và tính toán nằm ở backend.

\begin{figure}[H]
\centering
\resizebox{\textwidth}{!}{%
\begin{tikzpicture}[>=Latex, font=\small]
    \node[draw, rounded corners, minimum width=2.4cm, minimum height=0.8cm, align=center] (sv) at (0,0) {Sinh viên};
    \node[draw, rounded corners, minimum width=3.0cm, minimum height=0.8cm, align=center] (chatui) at (3.4,0) {ChatBot UI};
    \node[draw, rounded corners, minimum width=3.4cm, minimum height=0.8cm, align=center] (route) at (7.2,0) {\texttt{/chatbot/chat-stream}};
    \node[draw, rounded corners, minimum width=3.6cm, minimum height=0.8cm, align=center] (service) at (11.5,0) {\texttt{ChatbotService}\\và rule engine};
    \node[draw, rounded corners, minimum width=3.4cm, minimum height=0.8cm, align=center] (data) at (15.7,0) {Dữ liệu học tập\\\texttt{Student}, \texttt{LearnedSubject}, \texttt{CourseSubject}, \texttt{Subject}};

    \draw[dashed] (0,-0.6) -- (0,-6.5);
    \draw[dashed] (3.4,-0.6) -- (3.4,-6.5);
    \draw[dashed] (7.2,-0.6) -- (7.2,-6.5);
    \draw[dashed] (11.5,-0.6) -- (11.5,-6.5);
    \draw[dashed] (15.7,-0.6) -- (15.7,-6.5);

    \draw[->] (0.7,-1.1) -- node[above]{Nhập câu hỏi gợi ý học phần} (2.9,-1.1);
    \draw[->] (2.9,-2.0) -- node[above]{Gửi streaming request} (7.0,-2.0);
    \draw[->] (7.0,-3.0) -- node[above]{Phân loại intent\\\texttt{subject\_registration\_suggestion}} (11.2,-3.0);
    \draw[->] (11.2,-4.0) -- node[above]{Đọc CPA, học phần đã học\\và học phần trong chương trình} (15.4,-4.0);
    \draw[->] (15.4,-5.0) -- node[above]{Sinh danh sách học phần nên đăng ký} (11.7,-5.0);
    \draw[->] (11.7,-5.9) -- node[above]{Trả data và text cho FE\\để render bảng gợi ý} (3.7,-5.9);
\end{tikzpicture}
}%
\caption{Luồng truyền thông điệp khi sinh viên hỏi nên đăng ký học phần nào}
\end{figure}

Biểu đồ trên mô tả đúng cách FE hiện tại làm việc với chatbot: người dùng nhập câu hỏi, giao diện đẩy câu hỏi lên API streaming, backend phân loại intent rồi gọi phần tư vấn học phần. Kết quả trả về không chỉ là văn bản mà còn có danh sách học phần để giao diện hiển thị dưới dạng bảng hoặc danh sách gợi ý.

\begin{figure}[H]
\centering
\resizebox{\textwidth}{!}{%
\begin{tikzpicture}[>=Latex, font=\small]
    \node[draw, rounded corners, minimum width=2.4cm, minimum height=0.8cm, align=center] (sv) at (0,0) {Sinh viên};
    \node[draw, rounded corners, minimum width=3.0cm, minimum height=0.8cm, align=center] (chatui) at (3.4,0) {ChatBot UI};
    \node[draw, rounded corners, minimum width=3.4cm, minimum height=0.8cm, align=center] (route) at (7.2,0) {\texttt{/chatbot/chat-stream}};
    \node[draw, rounded corners, minimum width=3.7cm, minimum height=0.8cm, align=center] (service) at (11.6,0) {\texttt{ChatbotService}\\điều phối hội thoại};
    \node[draw, rounded corners, minimum width=3.8cm, minimum height=0.8cm, align=center] (query) at (16.1,0) {\texttt{ClassQueryService}\\và \texttt{ScheduleCombinationService}};
    \node[draw, rounded corners, minimum width=3.5cm, minimum height=0.8cm, align=center] (db) at (20.6,0) {\texttt{Class}, \texttt{Subject}, \texttt{ClassRegister}};

    \draw[dashed] (0,-0.6) -- (0,-6.9);
    \draw[dashed] (3.4,-0.6) -- (3.4,-6.9);
    \draw[dashed] (7.2,-0.6) -- (7.2,-6.9);
    \draw[dashed] (11.6,-0.6) -- (11.6,-6.9);
    \draw[dashed] (16.1,-0.6) -- (16.1,-6.9);
    \draw[dashed] (20.6,-0.6) -- (20.6,-6.9);

    \draw[->] (0.7,-1.1) -- node[above]{Nhập câu hỏi gợi ý lớp học} (2.9,-1.1);
    \draw[->] (2.9,-2.0) -- node[above]{Gửi streaming request} (7.0,-2.0);
    \draw[->] (7.0,-3.0) -- node[above]{Phân loại intent\\\texttt{class\_registration\_suggestion}} (11.3,-3.0);
    \draw[->] (11.3,-4.0) -- node[above]{Hỏi thêm preference nếu chưa đủ\\và lưu trạng thái hội thoại} (15.8,-4.0);
    \draw[->] (15.8,-5.0) -- node[above]{Lọc lớp theo ngày, giờ, giảng viên\\và tạo phương án lịch học} (20.3,-5.0);
    \draw[->] (20.3,-6.0) -- node[above]{Trả các phương án cho FE\\và cho phép xuất Excel} (3.8,-6.0);
\end{tikzpicture}
}%
\caption{Luồng truyền thông điệp khi sinh viên hỏi nên đăng ký lớp nào}
\end{figure}

Điểm đáng chú ý của luồng này là chatbot không trả lời một lần duy nhất mà có thể hỏi thêm thông tin nếu còn thiếu ràng buộc. Sau khi đủ preference, backend mới sinh các phương án lớp học và frontend hiển thị thành bảng lịch học trực quan, kèm tùy chọn xuất file Excel cho phương án phù hợp.

\section{Phát triển frontend}
Frontend được tổ chức theo hai layout chính: layout sinh viên và layout quản trị viên. Cách chia này giúp mỗi nhóm người dùng chỉ nhìn thấy đúng các chức năng phù hợp với vai trò của mình.

\subsection{Khu vực sinh viên}
\begin{itemize}
    \item \textbf{Dashboard}: hiển thị CPA hiện tại, số tín chỉ, số môn đã hoàn thành, mức cảnh báo học tập, điểm gần đây và lịch học sắp tới.
    \item \textbf{Lịch học và điểm số}: cho phép xem lịch học, danh sách môn đã học, GPA/CPA và các thống kê liên quan.
    \item \textbf{Đăng ký học phần và lớp học}: hỗ trợ tìm kiếm, lọc và gửi yêu cầu đăng ký theo từng kỳ.
    \item \textbf{Biểu mẫu, FAQ, phản hồi}: cung cấp kênh liên hệ và gửi yêu cầu hỗ trợ.
\end{itemize}

\subsection{Khu vực quản trị}
\begin{itemize}
    \item \textbf{StudentsManagement}: quản lý thông tin sinh viên.
    \item \textbf{SubjectsManagement}: quản lý học phần.
    \item \textbf{ClassesManagement, ScheduleManagement}: quản lý lớp học và thời khóa biểu.
    \item \textbf{FeedbackManagement, FAQManagement, NotificationManagement}: quản lý nội dung hỗ trợ và thông báo.
\end{itemize}

\subsection{Thành phần dùng chung}
\begin{itemize}
    \item \texttt{ProtectedRoute} kiểm tra quyền truy cập theo vai trò.
    \item \texttt{AuthContext} quản lý trạng thái đăng nhập.
    \item \texttt{ChatBot} là thành phần hội thoại trung tâm của hệ thống.
    \item Các component hỗ trợ như upload Excel, modal biểu mẫu và preview tài liệu giúp tái sử dụng giao diện tốt hơn.
\end{itemize}

\section{Phát triển chatbot}
\subsection{Mục tiêu}
Chatbot được xây dựng để hỗ trợ sinh viên tra cứu thông tin học tập, nhận gợi ý học phần/lớp học và lưu lại toàn bộ lịch sử hội thoại. Thay vì chỉ trả lời dạng văn bản cứng, chatbot có thể trả về bảng dữ liệu, danh sách lựa chọn hoặc phương án thời khóa biểu.

\subsection{Luồng xử lý}
\begin{enumerate}
    \item Người dùng nhập câu hỏi vào giao diện chat.
    \item Hệ thống chuẩn hóa câu hỏi và phân loại ý định.
    \item Tùy ý định, backend chuyển sang nhánh NL2SQL, rule engine hoặc agent orchestration.
    \item Kết quả được định dạng lại trước khi trả về frontend.
    \item Lịch sử cuộc trò chuyện được lưu vào CSDL để xem lại ở lần sau.
\end{enumerate}

\subsection{Các hướng xử lý chính}
\begin{itemize}
    \item \textbf{Tra cứu thông tin}: xem điểm, xem học phần, xem lớp, xem lịch học.
    \item \textbf{Gợi ý học phần}: dựa trên kết quả học tập và lộ trình chương trình đào tạo.
    \item \textbf{Gợi ý lớp học}: dựa trên sở thích về ngày học, giờ học, giáo viên và các ràng buộc xung đột lịch.
    \item \textbf{Xuất lịch học}: hỗ trợ tạo file Excel cho các phương án thời khóa biểu.
\end{itemize}

\subsection{Kết quả tích hợp}
Chatbot trên hệ thống không chỉ phục vụ hỏi đáp đơn giản mà còn đóng vai trò như một công cụ tư vấn học tập. Người dùng có thể nhập các câu hỏi gần với ngôn ngữ tự nhiên, hệ thống sẽ cố gắng hiểu ngữ cảnh, lọc dữ liệu phù hợp và trả về đáp án dễ đọc.

\subsection{Mã nguồn triển khai}
Các thành phần chính của chatbot được chia thành:
\begin{itemize}
    \item \texttt{tfidf\_classifier.py}: phân loại ý định.
    \item \texttt{chatbot\_service.py}: điều phối nghiệp vụ trả lời.
    \item \texttt{nl2sql\_service.py}: tạo truy vấn dữ liệu từ ngôn ngữ tự nhiên.
    \item \texttt{class\_query\_service.py}: truy vấn lớp học theo ràng buộc an toàn.
    \item \texttt{schedule\_combination\_service.py}: sinh tổ hợp phương án thời khóa biểu.
\end{itemize}

\section{Kết quả hoạt động}
Sau khi hoàn thành các phần lõi, hệ thống đạt được một số kết quả chính sau:
\begin{itemize}
    \item Người dùng có thể đăng ký, đăng nhập và khôi phục mật khẩu.
    \item Sinh viên xem được dashboard cá nhân, lịch học, điểm số, chương trình đào tạo và các mục hỗ trợ.
    \item Quản trị viên thao tác được với sinh viên, học phần, lớp học, thời khóa biểu, FAQ và thông báo.
    \item Chatbot hỗ trợ hỏi đáp học tập, tra cứu dữ liệu và gợi ý đăng ký học phần/lớp học.
    \item Lịch sử hội thoại, xuất Excel và các tiện ích phụ trợ hoạt động đồng bộ với backend.
\end{itemize}

Để thuận tiện trình bày trong báo cáo, có thể mô tả các nhóm màn hình theo bảng sau.

\begin{table}[H]
\centering
\begin{tabularx}{\textwidth}{|l|X|}
\hline
\textbf{Nhóm màn hình} & \textbf{Mô tả ngắn} \\ \hline
\textbf{Xác thực} & Đăng nhập, đăng ký, quên mật khẩu, đặt lại mật khẩu. \\ \hline
\textbf{Sinh viên} & Dashboard, lịch học, điểm số, chương trình đào tạo, đăng ký học phần, đăng ký lớp, biểu mẫu, FAQ, phản hồi. \\ \hline
\textbf{Quản trị} & Quản lý sinh viên, học phần, lớp học, thời khóa biểu, phản hồi, FAQ và thông báo. \\ \hline
\textbf{Chatbot} & Hỏi đáp học tập, gợi ý học phần, gợi ý lớp học, lưu lịch sử và xuất phương án lịch học. \\ \hline
\end{tabularx}
\caption{Các nhóm màn hình chính của hệ thống}
\end{table}

\section{Kết luận và hướng phát triển}
\subsection{Kết luận}
Đồ án đã xây dựng được một hệ thống quản lý sinh viên có tích hợp chatbot hỗ trợ học tập, cho phép tra cứu dữ liệu học tập, quản lý lớp học và hỗ trợ đăng ký học phần/lớp học theo hướng tự động hơn. Hệ thống có cấu trúc rõ ràng giữa frontend, backend và CSDL, thuận lợi cho việc mở rộng trong các giai đoạn sau.

\subsection{Hướng phát triển}
\begin{itemize}
    \item Nâng cấp chatbot bằng mô hình ngôn ngữ lớn để hiểu tốt hơn các câu hỏi tự nhiên và tình huống phức tạp.
    \item Hoàn thiện thêm các cơ chế đề xuất thời khóa biểu tối ưu hơn theo nhiều tiêu chí cùng lúc.
    \item Bổ sung thông báo thời gian thực, hệ thống cảnh báo học tập và đồng bộ lịch học tự động.
    \item Tăng cường kiểm thử tự động và hoàn thiện bộ dữ liệu mẫu để hỗ trợ triển khai thực tế.
\end{itemize}
\end{document}
